\documentclass[11pt]{article}
\usepackage[margin=2.8cm]{geometry}
\usepackage{amsmath,amssymb,amsthm}
\newtheorem{theorem}{Theorem}
\newtheorem{lemma}[theorem]{Lemma}
\newtheorem{corollary}[theorem]{Corollary}
\newtheorem{remark}[theorem]{Remark}
\newcommand{\qbin}[2]{\genfrac{[}{]}{0pt}{}{#1}{#2}}
\newcommand{\poch}[1]{(q;q)_{#1}}
\title{Warnaar's cylindric conjecture for the balanced profile $(4,3,3)$
at modulus $13$:\\ a complete proof via Warnaar's finite form ($a=+1$)
and Uncu's modulo-13 theorem}
\author{Seed 1, Layer 4, Round 2 (loop experiment 2026-07-04)}
\date{2026-07-04}
\begin{document}
\maketitle

\begin{abstract}
We prove the case $k=4$ of Warnaar's Conjecture
\texttt{Con\_cylindric-b} (the $\mathrm{A}_2$ Andrews--Gordon
fermionic-form conjecture for cylindric partitions of balanced profile
$(k,k-1,k-1)$ at modulus $3k+1$): the profile $c=(4,3,3)$ at level
$d=10$ (modulus $13$). The proof chains two proved results:
Warnaar's finite-form Proposition ($a=+1$, $k=3$) in a coefficientwise
limit, and Uncu's 2024 modulo-13 theorem expressing $H_{(4,3,3)}$ in the
series $S_{13}(\rho|\sigma)$. Unlike the $d=8$ balanced case, \emph{no}
Kanade--Russell contiguous relation is needed: after the Pochhammer
split, the limit of the finite form coincides verbatim with Uncu's proved
expression. As a corollary the coefficients $Q_{n,(4,3,3)}(q)$ are
manifestly nonnegative: the second proved balanced case of the
fermionic-form conjecture, the first proved case of
\texttt{Con\_cylindric-b} beyond $k=2$, and the first proved-positive
core orbit at $d=10$.
\end{abstract}

\section{Setup and statement}

Throughout, $\poch{n}=\prod_{i=1}^n(1-q^i)$, and $\qbin{n}{m}$ is the
$q$-binomial coefficient (zero unless $0\le m\le n$). For a profile
$c=(c_1,c_2,c_3)$ with $c_1+c_2+c_3=d$, let $F_c(z,q)$ be the
Corteel--Welsh cylindric-partition generating function
(raw interlacing definition; $z$ tracks the largest part, $q$ the size),
and set
\[
G_c(z,q)=(zq;q)_\infty F_c(z,q),\qquad
H_c(z,q)=\frac{(zq;q)_\infty}{(q;q)_\infty}F_c(z,q)=\frac{G_c(z,q)}{(q;q)_\infty},
\]
\[
Q_{n,c}(q)=\poch{n}\,[z^n]\,G_c(z,q).
\]
Warnaar's positivity conjecture (Corteel--Dousse--Uncu) asserts
$Q_{n,c}(q)\ge 0$ coefficientwise for all $n,c$ with
$d\not\equiv0\pmod 3$. For the balanced profile $(k,k-1,k-1)$ of
modulus $3k+1$, Warnaar's Conjecture \texttt{Con\_cylindric-b}
(first identity) predicts the fermionic form
\begin{equation}\label{eq:ferm}
G_{(k,k-1,k-1)}(z,q)\;\overset{?}{=}\;
\mathrm{FERM}^{+}_{k-1}(z,q):=
\sum_{\substack{n_1,\dots,n_{k-1}\ge 0\\ m_1,\dots,m_{k-1}\ge 0}}
\frac{z^{n_1}q^{\sum_{i=1}^{k-1}(n_i^2-n_im_i+m_i^2)}}{\poch{n_1}}\,
\qbin{2n_{k-1}}{m_{k-1}}
\prod_{i=1}^{k-2}
\qbin{n_i}{n_{i+1}}\qbin{n_i-n_{i+1}+m_{i+1}}{m_i}.
\end{equation}
The case $k=1$ is trivial and $k=2$ is Corteel--Welsh
\cite[Theorem~3.2]{CW19}; all $k\ge3$ were open. We prove the case
$k=4$, profile $(4,3,3)$, $d=10$, modulus $13$.

For $m=13$ ($k=4$ in Uncu's indexing $m=3k+1$), $\rho,\sigma\in\mathbb{Z}^3$,
define (Uncu 2024, eq.\ (\texttt{eq:Sp1}); note there is \emph{no}
$q^{2r_3s_3}$ cross term, in contrast to $S_{3k-1}$)
\begin{equation}\label{eq:S13}
S_{13}(z;\rho|\sigma)=
\sum_{\substack{r_1\ge r_2\ge r_3\ge 0\\ s_1\ge s_2\ge s_3\ge 0}}
z^{r_1}\,
\frac{q^{\sum_{i=1}^{3}(r_i^2-r_is_i+s_i^2+\rho_ir_i+\sigma_is_i)}}
{\poch{r_1-r_2}\poch{r_2-r_3}\poch{r_3}\,
 \poch{s_1-s_2}\poch{s_2-s_3}\poch{s_3}\,\poch{r_3+s_3+1}},
\end{equation}
and $e_3=(0,0,0)$, $e_2=(0,0,1)$. We suppress $z$ and write
$S(\rho|\sigma)$.

\begin{theorem}\label{thm:main}
$G_{(4,3,3)}(z,q)=\mathrm{FERM}^{+}_{3}(z,q)$, i.e.\ Warnaar's
Conjecture \emph{\texttt{Con\_cylindric-b}} (first identity) holds for
$k=4$.
\end{theorem}

\begin{corollary}\label{cor:pos}
For all $n\ge 0$,
\[
Q_{n,(4,3,3)}(q)=
\sum_{\substack{n_2,n_3\ge 0\\ m_1,m_2,m_3\ge 0}}
q^{\,n^2+n_2^2+n_3^2-nm_1-n_2m_2-n_3m_3+m_1^2+m_2^2+m_3^2}\,
\qbin{n}{n_2}\qbin{n-n_2+m_2}{m_1}
\qbin{n_2}{n_3}\qbin{n_2-n_3+m_3}{m_2}\qbin{2n_3}{m_3},
\]
a manifestly nonnegative polynomial (each $n_i^2-n_im_i+m_i^2\ge0$, all
$q$-binomials nonnegative). In particular Warnaar's positivity
conjecture holds for the profile $(4,3,3)$ at $d=10$.
\end{corollary}

\begin{proof}[Proof of the corollary]
$Q_{n,(4,3,3)}=\poch{n}[z^n]\mathrm{FERM}^{+}_3$; the summand of
\eqref{eq:ferm} at $n_1=n$ has denominator $\poch{n_1}=\poch{n}$, which
cancels.
\end{proof}

\begin{remark}
The profile $(4,3,3)$ is chirality-safe: its reversal $(3,3,4)$ is a
cyclic shift of $(4,3,3)$, and $F_c$ is invariant under cyclic shifts of
the profile. No orbit-labelling subtleties (synthesis-layer3 \S4(iv))
arise.
\end{remark}

\section{Proof of Theorem \ref{thm:main}}

The proof chains three proved statements. In contrast to the $d=8$
balanced case (Seed 2, Layer 3), no bridging contiguous relation is
required: the modulus-$13$ family has no $q^{2r_3s_3}$ cross term, and
the Pochhammer split lands directly on Uncu's proved expression.

\subsection*{Step 1: Warnaar's finite form ($a=+1$) in the limit}
Warnaar ($\mathrm{A}_2$ Andrews--Gordon paper, Section~7) defines, for
$\sigma=(\sigma_1,\dots,\sigma_k)=(1,\dots,1,a)$, $a\in\{-1,1\}$,
\[
F^{(a)}_{n_0,m_0;k}(z,q)=
\sum_{\substack{n_1,\dots,n_k\ge 0\\ m_1,\dots,m_k\ge 0}}
\frac{z^{n_1}}{\poch{n_k+m_k}}
\prod_{i=1}^{k}q^{n_i^2-\sigma_in_im_i+m_i^2}
\qbin{n_{i-1}}{n_i}\qbin{m_{i-1}}{m_i},
\]
and proves (his Proposition \texttt{Prop\_finiteform}, case $a=+1$, in
which $\sigma_i=1$ for all $i$) the identity
\begin{equation}\label{eq:ff}
F^{(1)}_{n_0,m_0;k}(z,q)=
\sum_{\substack{n_1,\dots,n_k\ge 0\\ m_1,\dots,m_k\ge 0}}
\frac{z^{n_1}q^{\sum_{i=1}^{k}(n_i^2-n_im_i+m_i^2)}}{\poch{m_0-n_1+m_1}}\,
\qbin{n_0}{n_1}\qbin{2n_k}{m_k}
\prod_{i=1}^{k-1}
\qbin{n_i}{n_{i+1}}\qbin{n_i-n_{i+1}+m_{i+1}}{m_i}.
\end{equation}
Fix $k=3$ and let $n_0,m_0\to\infty$. Each coefficient of $z^{n}q^{N}$ on
both sides stabilizes: on the right, $\qbin{n_0}{n_1}\to 1/\poch{n_1}$
and $1/\poch{m_0-n_1+m_1}\to 1/\poch{\infty}$ coefficientwise; on the
left the binomial chains telescope,
$\qbin{n_0}{n_1}\qbin{n_1}{n_2}\qbin{n_2}{n_3}\to
1/(\poch{n_1-n_2}\poch{n_2-n_3}\poch{n_3})$ and likewise for the
$m$-chain. (Stabilization: for a fixed monomial $z^nq^N$, all summands
with any index exceeding $N+n$ contribute only above $q^N$, so both
sides are finite sums of eventually constant coefficients.)
Writing $(r,s)=(n,m)$, the left side becomes
\begin{equation}\label{eq:T}
\frac{\mathrm{FERM}^{+}_3(z,q)}{(q;q)_\infty}=T(z,q):=
\sum_{\substack{r_1\ge r_2\ge r_3\ge 0\\ s_1\ge s_2\ge s_3\ge 0}}
z^{r_1}\,
\frac{q^{\sum_{i=1}^{3}(r_i^2-r_is_i+s_i^2)}}
{\poch{r_1-r_2}\poch{r_2-r_3}\poch{r_3}\,
 \poch{s_1-s_2}\poch{s_2-s_3}\poch{s_3}\,\poch{r_3+s_3}},
\end{equation}
where the right side of \eqref{eq:ff} has become
$\mathrm{FERM}^{+}_3(z,q)/(q;q)_\infty$.

\subsection*{Step 2: Pochhammer split}
Using $\dfrac{1}{\poch{r_3+s_3}}=\dfrac{1-q^{r_3+s_3+1}}{\poch{r_3+s_3+1}}$
and $q^{r_3+s_3+1}=q\cdot q^{r_3}\cdot q^{s_3}$ (a shift
$(\rho_3,\sigma_3)\mapsto(\rho_3+1,\sigma_3+1)$ in \eqref{eq:S13}),
term by term,
\begin{equation}\label{eq:split}
T(z,q)=S(e_3|e_3)-q\,S(e_2|e_2)
     =S_{13}\bigl((0,0,0)|(0,0,0)\bigr)-q\,S_{13}\bigl((0,0,1)|(0,0,1)\bigr).
\end{equation}

\subsection*{Step 3: Uncu's modulo-13 theorem}
Uncu (2024, Theorem \texttt{thm:m13}, first proved via Theorem
\texttt{thm:n6m13}) proves the Kanade--Russell expression
(\texttt{eq:mod13list}, last row) for the profile $(4,3,3)$:
\begin{equation}\label{eq:uncu}
H_{(4,3,3)}(z,q)=
S_{13}\bigl((0,0,0)|(0,0,0)\bigr)-q\,S_{13}\bigl((0,0,1)|(0,0,1)\bigr).
\end{equation}
Uncu's normalization $H_c=(zq;q)_\infty F_c/(q;q)_\infty$
(his eq.\ (\texttt{eq:HtoF})) agrees with ours, and his $F_c$ is the raw
interlacing generating function.

\subsection*{Conclusion}
Combining \eqref{eq:T}--\eqref{eq:uncu},
\[
\frac{\mathrm{FERM}^{+}_3(z,q)}{(q;q)_\infty}=T(z,q)=H_{(4,3,3)}(z,q)
=\frac{G_{(4,3,3)}(z,q)}{(q;q)_\infty},
\]
hence $G_{(4,3,3)}=\mathrm{FERM}^{+}_3$. \qed

\section{Numerical verification and remarks}

\begin{remark}[Verification artifacts]
All steps were verified in exact arithmetic
(\texttt{scratch/scripts/seed1\_R2L4\_d10\_chain.sage}, SageMath 10.9;
log \texttt{scratch/tmp/seed1\_L4\_d10\_chain\_n12.log}):
\begin{itemize}
\item[{[FF]}] Warnaar's finite form \eqref{eq:ff} ($a=+1$, $k=3$) as an
exact power-series identity (to $q^{150}$, computed at $q^{300}$) for all
$0\le n_0,m_0\le 3$ --- a transcription guard on the quoted proposition.
\item[{[A]}] The split \eqref{eq:split} at $z$-orders $n=0,\dots,4$ to
$q^{150}$.
\item[{[B]}] $\poch{n}(q;q)_\infty[z^n]\bigl(S(e_3|e_3)-qS(e_2|e_2)\bigr)
=Q_{n,(4,3,3)}$ against the reference engine
(\texttt{seed8\_R2L3\_engine.sage} recursion: target-first kernel, exact
$\mathbb{Z}[q]$, exact $\Phi_3(q^m)$ division, $Q_n$ by Gauss inversion,
Theorem~Q/Corollary~I of Seed 7 Layer 3) for $n=0,\dots,4$ to $q^{150}$.
\item[{[C]}] The finite fermionic sum of Corollary~\ref{cor:pos} equals
the engine $Q_{n,(4,3,3)}$ \emph{exactly in} $\mathbb{Z}[q]$ for
$n=0,\dots,12$ ($\deg Q_{12}=1296$), with all coefficients nonnegative;
also $Q_n(1)=21^n$ ($K-1=21$ at $d=10$), as required.
\item[{[D]}] The coefficientwise limit: $\poch{n}(q;q)_\infty[z^n]T$
equals the fermionic sum for $n=0,\dots,4$ to $q^{150}$.
\end{itemize}
\end{remark}

\begin{remark}[On the Kanade--Russell relation and Uncu's $R_3$ typo]
The Layer-3 brief cautioned that Uncu's displayed $R_3$ relation for
$m\equiv1\pmod 3$ contains a typo (two identical terms of opposite sign;
his eq.\ (\texttt{R31})). This is confirmed --- the displayed relation is
vacuous as printed --- but it is \emph{irrelevant to this proof}: at
modulus $13$ the Pochhammer split \eqref{eq:split} produces exactly the
pair $(e_3|e_3)$, $(e_2|e_2)$ appearing in Uncu's proved formula
\eqref{eq:uncu}, so no contiguous relation is needed. Structurally: at
$m=3k-1$ (the $d=8$ case) the $a=-1$ finite form produces the
$(e_3|e_3),(e_2|e_2)$ pair while the Kanade--Russell table lists
$(e_3|e_2),(e_2|e_1)$, requiring one $R_3$; at $m=3k+1$ the
Kanade--Russell pattern $H_{(c_1,c_2,c_3)}=S(e_{c_2}|e_{c_3})
-qS(e_{c_2-1}|e_{c_3-1})$ evaluated at $c_2=c_3=k-1$ gives exactly the
split pair.
\end{remark}

\begin{remark}[Extensions]
(i) For general $k$, the identical chain reduces
\texttt{Con\_cylindric-b} (first identity) at modulus $3k+1$ to the
Kanade--Russell/Uncu expression for $H_{(k,k-1,k-1)}$, proved to date
only for $m=13$. The $k=3$ case ($d=7$, modulus $10$) remains the
smallest balanced case not proved anywhere.
(ii) With Seed 2's $d=8$ result, both balanced-orbit templates
($a=\pm1$) are now realized; the remaining $d=10$ core orbits require
the $R_1/R_2$ word search plus new finite-form seeds, exactly as at
$d=8$ (synthesis-layer3 \S4(v)).
\end{remark}

\begin{thebibliography}{9}
\bibitem{CW19} S.~Corteel, T.~Welsh,
\emph{The $\mathrm{A}_2$ Rogers--Ramanujan identities revisited},
Ann. Comb. 23 (2019).
\bibitem{W23} S.~O.~Warnaar,
\emph{The $\mathrm{A}_2$ Andrews--Gordon identities and cylindric
partitions}, Trans. Amer. Math. Soc. Ser. B (2023).
(Local: \texttt{literature/tex/warnaar\_A2\_andrews\_gordon/source.tex};
Proposition \texttt{Prop\_finiteform} lines 2672--2687, proof lines
2784--2819; Conjecture \texttt{Con\_cylindric-b} lines 767--785.)
\bibitem{U24} A.~K.~Uncu,
\emph{Proofs of modulo 11 and 13 cylindric Kanade--Russell conjectures
for $\mathrm{A}_2$ Rogers--Ramanujan type identities}, 2024.
(Local: \texttt{literature/corteel-citations/tex/uncu\_proofs\_modulo11\_13\_cylindric\_kanade\_russell/main.tex};
$S_{13}$ eq.\ (\texttt{eq:Sp1}) line 235; \texttt{eq:mod13list} lines
452--483; Theorem \texttt{thm:m13} line 573.)
\bibitem{KR22} S.~Kanade, M.~C.~Russell,
\emph{Completing the $\mathrm{A}_2$ Andrews--Schilling--Warnaar
identities}, 2022. (Contiguous relations; not needed in this proof.)
\end{thebibliography}

\end{document}
